% Differentiable Operator Calibration for Coded Aperture
% Snapshot Spectral Imaging
%
% LaTeX source - ECCV/LNCS format

\documentclass[runningheads]{llncs}
\usepackage[T1]{fontenc}
\usepackage{graphicx}
\usepackage{amsmath,amssymb}
\usepackage{booktabs}
\usepackage{multirow}
\usepackage{xcolor}
\usepackage{hyperref}
\usepackage{cleveref}
\usepackage{subcaption}
\usepackage{array}
\usepackage{algorithm}
\usepackage{algorithmic}
\usepackage{tikz}
\usetikzlibrary{shapes,arrows,positioning,fit,calc}

% Custom column type for tables
\newcolumntype{C}[1]{>{\centering\arraybackslash}p{#1}}

% Color definitions
\definecolor{bestcolor}{rgb}{0.0, 0.5, 0.0}
\definecolor{worstcolor}{rgb}{0.7, 0.0, 0.0}
\definecolor{stagecolor}{rgb}{0.2, 0.4, 0.8}

\newcommand{\best}[1]{\textcolor{bestcolor}{\textbf{#1}}}
\newcommand{\worst}[1]{\textcolor{worstcolor}{#1}}

\begin{document}

\title{Differentiable Operator Calibration for Coded Aperture Snapshot Spectral Imaging}
\titlerunning{Differentiable CASSI Calibration}

\author{Chengshuai Yang}
\authorrunning{C. Yang}
\institute{NextGen PlatformAI C Corp}

\maketitle

% ============================================================================
% ABSTRACT
% ============================================================================
\begin{abstract}
Coded aperture snapshot spectral imaging (CASSI) acquires hyperspectral data cubes in a single shot but requires accurate knowledge of the forward measurement operator---the coded aperture mask position, orientation, and dispersive element parameters---for high-quality reconstruction.
In practice, manufacturing tolerances and assembly drift introduce operator mismatch that degrades reconstruction by 10--17\,dB.
We present a differentiable calibration framework that models CASSI mismatch as a 6-parameter perturbation (spatial shift, rotation, dispersion slope and axis angle) and recovers these parameters through a two-stage pipeline:
(1)~a hierarchical beam search over a coarse parameter grid ($\sim$38\,s/scene), followed by
(2)~a joint gradient refinement using differentiable PyTorch modules---including a straight-through estimator for integer dispersion offsets and an unrolled GAP-TV solver with gradient checkpointing ($\sim$366\,s/scene).
Central to our approach is an enlarged grid forward model with $4\times$ spatial and $2\times$ spectral oversampling (217 bands), providing sub-pixel sensitivity to mismatch parameters.
Validated on 10 KAIST hyperspectral scenes under a three-scenario protocol, our method achieves a calibration gain of $+5.06$\,dB, recovering 30\% of the 16.60\,dB mismatch loss.
When combined with oracle correction using mask-aware deep networks (MST-L), the recovery reaches $+7.99$\,dB (75.5\% of mismatch loss), demonstrating the synergy between calibration and learned reconstruction.

\keywords{CASSI \and Operator calibration \and Differentiable programming \and Hyperspectral imaging \and Mismatch correction \and Computational imaging}
\end{abstract}

% ============================================================================
% 1. INTRODUCTION
% ============================================================================
\section{Introduction}
\label{sec:intro}

Coded aperture snapshot spectral imaging (CASSI) acquires a three-dimensional hyperspectral data cube $\mathbf{x} \in \mathbb{R}^{H \times W \times \Lambda}$ from a single two-dimensional measurement $\mathbf{y} \in \mathbb{R}^{H \times W'}$ by encoding spectral information through a binary coded aperture followed by a dispersive prism~\cite{wagadarikar2008cassi,gehm2007cassi}.
Computational reconstruction algorithms then recover the hyperspectral cube from this compressed measurement, exploiting sparsity in appropriate transform domains~\cite{yuan2016gaptv,cai2022mst}.

The quality of CASSI reconstruction depends critically on accurate knowledge of the forward measurement operator $\mathbf{\Phi}$---the mapping from scene to measurement defined by the coded aperture mask position, orientation, and dispersive element parameters.
In practice, the assumed operator $\hat{\mathbf{\Phi}}$ inevitably differs from the true physical operator $\mathbf{\Phi}$ due to manufacturing tolerances, mechanical assembly errors, thermal drift, and optical alignment imprecision.
Even moderate misalignment---0.5 pixels of mask shift and 0.1$^\circ$ of rotation---can degrade peak signal-to-noise ratio (PSNR) by over 10\,dB~\cite{wagadarikar2008cassi}, erasing the quality advantage of state-of-the-art deep learning reconstruction methods~\cite{cai2022mst,hu2022hdnet}.
As demonstrated by the InverseNet benchmark~\cite{yang2025inversenet}, operator mismatch reduces all deep learning methods to a narrow 24--25\,dB performance band regardless of their 10+ dB spread under ideal conditions.

Current CASSI calibration approaches are predominantly offline procedures requiring dedicated calibration targets (checkerboards, spectral lamps) and manual alignment~\cite{arguello2013cassi_calib,kittle2010cassi_design}.
These methods cannot adapt to in-situ drift and do not provide a differentiable framework suitable for gradient-based optimization.
More recently, differentiable programming has been applied to inverse problems in imaging~\cite{gilton2021deep_eq,monga2021unrolled}, but its application to CASSI operator calibration remains unexplored.

\textbf{Contributions.}
We address this gap with three contributions:
\begin{enumerate}
    \item \textbf{Enlarged grid forward model with 6-parameter mismatch.}
    We introduce a high-fidelity CASSI forward model with $N\!=\!4$ spatial oversampling and $K\!=\!2$ spectral oversampling (expanding 28 bands to 217), enabling sub-pixel sensitivity to a 6-parameter mismatch model covering spatial shift $(d_x, d_y)$, rotation $(\theta)$, dispersion slope $(a_1)$, and dispersion axis angle $(\alpha)$.

    \item \textbf{Two-stage calibration pipeline.}
    We propose a coarse-to-fine calibration strategy: Algorithm~1 performs hierarchical beam search over discrete parameter grids ($\sim$38\,s/scene), providing initial estimates that Algorithm~2 refines through joint gradient descent using differentiable PyTorch modules ($\sim$366\,s/scene), achieving 3--5$\times$ accuracy improvement.

    \item \textbf{Differentiable CASSI modules.}
    We develop four differentiable PyTorch modules---RoundSTE (straight-through estimator for integer offsets), DifferentiableMaskWarp (affine warping), DifferentiableCassiForward/Adjoint (forward model with STE), and DifferentiableGAPTV (unrolled solver with gradient checkpointing)---enabling end-to-end gradient-based calibration of the CASSI operator.
\end{enumerate}

% ============================================================================
% 2. RELATED WORK
% ============================================================================
\section{Related Work}
\label{sec:related}

\paragraph{CASSI systems and forward models.}
CASSI was introduced by Wagadarikar et al.~\cite{wagadarikar2008cassi} as a snapshot spectral imager using a binary coded aperture and dispersive prism.
Gehm et al.~\cite{gehm2007cassi} developed the dual-disperser (DD-CASSI) variant.
The standard forward model assumes a perfectly aligned mask with known dispersion, discretized at the detector pixel pitch.
Higher-order models incorporating sub-pixel effects were proposed by Arguello et al.~\cite{arguello2013cassi_calib}, but without differentiable implementations.
Kittle et al.~\cite{kittle2010cassi_design} analyzed the design space of multiframe CASSI systems.

\paragraph{Calibration in computational imaging.}
Calibration of compressive imaging systems has traditionally relied on offline procedures with dedicated targets~\cite{arguello2013cassi_calib}.
For MRI, coil sensitivity calibration~\cite{uecker2014espirit} and trajectory correction~\cite{deshmane2012parallel} are standard preprocessing steps.
In phase retrieval, model mismatch has been studied through alternating projections~\cite{elser2003phase}.
These approaches are modality-specific and do not provide a unified differentiable framework.

\paragraph{Differentiable programming for inverse problems.}
Deep unrolling~\cite{monga2021unrolled} replaces iterative algorithm steps with learnable modules, enabling end-to-end training.
Deep equilibrium models~\cite{gilton2021deep_eq} compute fixed points of implicit layers.
The straight-through estimator (STE)~\cite{bengio2013ste} enables gradient flow through discrete operations and has been widely adopted in quantization-aware training~\cite{jacob2018quantization}.
Our work applies STE to the integer dispersion offsets in CASSI, enabling gradient-based calibration of inherently discrete parameters.

\paragraph{CASSI reconstruction methods.}
Classical methods include GAP-TV~\cite{yuan2016gaptv} (generalized alternating projection with total variation) and TwIST~\cite{bioucas2007twist}.
Deep learning methods have achieved significant quality improvements: MST~\cite{cai2022mst} introduces mask-guided spectral transformers, HDNet~\cite{hu2022hdnet} uses dual-domain processing, and CST~\cite{cai2022cst} extends the transformer approach.
The InverseNet benchmark~\cite{yang2025inversenet} demonstrated that mask-aware methods (MST-L: 34.81\,dB ideal, 75.5\% oracle recovery) vastly outperform mask-oblivious methods (HDNet: 0\% recovery) when calibration is available, motivating the development of practical calibration algorithms.

% ============================================================================
% 3. CASSI FORWARD MODEL WITH MISMATCH
% ============================================================================
\section{CASSI Forward Model with Mismatch}
\label{sec:forward}

\subsection{Standard CASSI Forward Model}
\label{sec:standard}

The standard single-disperser CASSI acquires a 2D measurement $\mathbf{y} \in \mathbb{R}^{H \times W'}$ of a 3D hyperspectral cube $\mathbf{x} \in \mathbb{R}^{H \times W \times \Lambda}$ through a coded aperture mask $\mathbf{M} \in \{0,1\}^{H \times W}$ followed by a dispersive prism:
\begin{equation}
    y(i, j) = \sum_{\lambda=1}^{\Lambda} M(i, j - d(\lambda)) \cdot x(i, j, \lambda) + n(i, j),
    \label{eq:cassi_forward}
\end{equation}
where $d(\lambda) = s \cdot (\lambda - 1)$ is the dispersion shift for spectral band $\lambda$ with step size $s$ (typically 2 pixels/band), and $W' = W + (\Lambda - 1) \cdot s$ is the measurement width.
In matrix-vector form, $\mathbf{y} = \mathbf{\Phi}\mathbf{x} + \mathbf{n}$, where $\mathbf{\Phi}$ encodes both the coded aperture and the dispersion.

\subsection{Six-Parameter Mismatch Model}
\label{sec:mismatch}

We model CASSI operator mismatch as a 6-parameter perturbation $\boldsymbol{\xi} = (d_x, d_y, \theta, a_1, \alpha, \sigma)$ organized into three groups:

\paragraph{Group 1: Mask affine transform $(d_x, d_y, \theta)$.}
The coded aperture mask undergoes spatial shift $(d_x, d_y)$ and rotation $\theta$ due to mechanical assembly tolerances:
\begin{equation}
    \tilde{M}(i,j) = M\!\left(\mathbf{R}_{-\theta}\begin{bmatrix}i - c_y - d_y \\ j - c_x - d_x\end{bmatrix} + \begin{bmatrix}c_y \\ c_x\end{bmatrix}\right),
    \label{eq:mask_warp}
\end{equation}
where $\mathbf{R}_{-\theta}$ is the inverse rotation matrix and $(c_y, c_x)$ is the mask center.
Typical ranges: $d_x, d_y \in [-3, 3]$\,px, $\theta \in [-1^\circ, 1^\circ]$.

\paragraph{Group 2: Dispersion parameters $(a_1, \alpha)$.}
Thermal drift and prism settling cause the dispersion slope $a_1$ (pixels per spectral band) and axis angle $\alpha$ to deviate from nominal values:
\begin{equation}
    d(\lambda) = \text{round}(a_1 \cdot (\lambda - \lambda_c)), \quad \text{with axis rotation } \alpha,
    \label{eq:dispersion}
\end{equation}
where $\lambda_c$ is the center wavelength.
Typical ranges: $a_1 \in [1.95, 2.05]$, $\alpha \in [-1^\circ, 1^\circ]$.

\paragraph{Group 3: PSF blur $(\sigma)$.}
Optional Gaussian PSF blur from lens misalignment, with $\sigma \in [0.5, 2.0]$\,px.
We find this parameter has low impact ($<$0.1\,dB) and do not actively correct it.

The mismatched forward operator is then:
\begin{equation}
    \tilde{y}(i, j) = \sum_{\lambda=1}^{\Lambda} \tilde{M}(i, j - \tilde{d}(\lambda)) \cdot x(i, j, \lambda) + n(i, j),
\end{equation}
where $\tilde{M}$ is the warped mask (Eq.~\ref{eq:mask_warp}) and $\tilde{d}$ uses the perturbed dispersion (Eq.~\ref{eq:dispersion}).

\subsection{Enlarged Grid Forward Model}
\label{sec:enlarged}

To achieve sub-pixel sensitivity, we introduce an enlarged grid forward model with spatial oversampling factor $N\!=\!4$ and spectral oversampling factor $K\!=\!2$:
\begin{itemize}
    \item \textbf{Spatial enlargement:} The $H \times W$ scene and mask are upsampled to $NH \times NW$ (e.g., $256 \to 1024$).
    \item \textbf{Spectral expansion:} The $\Lambda$ bands are interpolated to $L_{\text{exp}} = (\Lambda-1) \cdot N \cdot K + 1$ bands (e.g., $28 \to 217$).
    \item \textbf{Stride-1 dispersion:} In the enlarged space, each spectral band shifts by 1~pixel (vs.\ $s\!=\!2$ at native resolution), enabling fine-grained dispersion modeling.
    \item \textbf{Downsampling:} The enlarged measurement $(NH \times NW')$ is block-averaged back to $(H \times W')$.
\end{itemize}

This model provides $4\times$ sub-pixel spatial resolution for mismatch parameter sensitivity.
The measurement width in enlarged space is $NW + (L_{\text{exp}}-1) = 1024 + 216 = 1240$.

\subsection{Three-Scenario Evaluation Protocol}
\label{sec:protocol}

Following the InverseNet protocol~\cite{yang2025inversenet}, we evaluate under three scenarios:
\begin{itemize}
    \item \textbf{Scenario~I (Ideal):} Measurement with ideal operator, reconstruction with ideal operator.
    \item \textbf{Scenario~II (Mismatched):} Measurement with mismatched operator, reconstruction with assumed (ideal) operator.
    \item \textbf{Scenario~III (Calibrated):} Measurement with mismatched operator, reconstruction with operator corrected by Algorithms~1+2.
\end{itemize}

The calibration gain is $\Delta_{\text{cal}} = \text{PSNR}_{\text{III}} - \text{PSNR}_{\text{II}}$, quantifying the benefit of our calibration pipeline.
The residual gap $\Delta_{\text{res}} = \text{PSNR}_{\text{I}} - \text{PSNR}_{\text{III}}$ measures unrecoverable losses.

% ============================================================================
% 4. CALIBRATION ALGORITHMS
% ============================================================================
\section{Calibration Algorithms}
\label{sec:algorithms}

\subsection{Algorithm 1: Hierarchical Beam Search}
\label{sec:alg1}

Algorithm~1 performs coarse parameter estimation through discrete grid search, exploiting the separability of the mismatch parameter groups.

\begin{algorithm}[t]
\caption{Hierarchical Beam Search (Coarse Calibration)}
\label{alg:beam}
\begin{algorithmic}[1]
\REQUIRE Measurement $\mathbf{y}$, nominal mask $\mathbf{M}$, solver $\mathcal{R}$
\ENSURE Estimated mismatch $\hat{\boldsymbol{\xi}}_1 = (\hat{d}_x, \hat{d}_y, \hat{\theta}, \hat{a}_1, \hat{\alpha})$
\STATE \textbf{Coarse 3D grid:} Evaluate $9 \times 9 \times 7 = 567$ candidates over $(d_x, d_y, \theta)$
\STATE \hspace{1em} $d_x \in \{-2, -1.5, \ldots, 2\}$, $d_y \in \{-2, -1.5, \ldots, 2\}$, $\theta \in \{-0.75, \ldots, 0.75\}$
\STATE Select top-$k$ candidates by reconstruction PSNR
\STATE \textbf{Fine 3D beam search:} $5 \times 5 \times 5$ refinement around each top candidate
\STATE \textbf{2D beam search:} $5 \times 7$ grid over dispersion $(a_1, \alpha)$
\STATE \hspace{1em} $a_1 \in \{1.95, \ldots, 2.05\}$, $\alpha \in \{-1, \ldots, 1\}^\circ$
\STATE \textbf{Coordinate descent:} 3 rounds of local optimization over all parameters
\RETURN Best $\hat{\boldsymbol{\xi}}_1$
\end{algorithmic}
\end{algorithm}

The coarse 3D grid evaluates 567 forward-model/reconstruction pairs, each using a short GAP-TV run (8--12 iterations) as the quality metric.
The hierarchical structure reduces the total evaluation count from an exhaustive $O(n^5)$ grid to approximately 700 evaluations.
Complexity: $\sim$38\,s/scene on a single CPU core.

\subsection{Algorithm 2: Joint Gradient Refinement}
\label{sec:alg2}

Algorithm~2 refines the coarse estimate from Algorithm~1 using gradient-based optimization with differentiable CASSI modules.
The pipeline consists of five stages with increasing parameter scope:

\begin{algorithm}[t]
\caption{Joint Gradient Refinement (Fine Calibration)}
\label{alg:gradient}
\begin{algorithmic}[1]
\REQUIRE Measurement $\mathbf{y}$, nominal mask $\mathbf{M}$, initial estimate $\hat{\boldsymbol{\xi}}_1$
\ENSURE Refined estimate $\hat{\boldsymbol{\xi}}_2$
\STATE \textbf{Stage 0 (Coarse GPU grid):} 567 GPU-parallel evaluations, 8-iter GAP-TV ($\sim$85\,s)
\STATE \textbf{Stage 1 (Fine GPU grid):} 375 evaluations around Stage~0 best, 12-iter GAP-TV ($\sim$88\,s)
\STATE \textbf{Stage 2A (Gradient $d_x$):} 50 Adam steps, $\text{lr}=10^{-2}$, smoothing $\sigma=0.5$ ($\sim$61\,s)
\STATE \textbf{Stage 2B (Gradient $d_y, \theta$):} 60 Adam steps, $\text{lr}=10^{-2}$, $\sigma=1.0$ ($\sim$74\,s)
\STATE \textbf{Stage 2C (Joint refinement):} 80 Adam steps on $(d_x, d_y, \theta)$, $\sigma=0.7$ ($\sim$128\,s)
\RETURN $\hat{\boldsymbol{\xi}}_2$
\end{algorithmic}
\end{algorithm}

Each gradient step minimizes the reconstruction loss:
\begin{equation}
    \mathcal{L}(\boldsymbol{\xi}) = \|\mathbf{y} - \mathbf{\Phi}(\boldsymbol{\xi})\,\mathcal{R}(\mathbf{y}, \mathbf{\Phi}(\boldsymbol{\xi}))\|_2^2,
    \label{eq:loss}
\end{equation}
where $\mathcal{R}(\mathbf{y}, \mathbf{\Phi})$ is the unrolled GAP-TV reconstruction.
The staged approach addresses parameter coupling: $d_x$ is refined first (Stage~2A) because it has the strongest gradient signal, followed by the coupled pair $(d_y, \theta)$ (Stage~2B), then joint refinement (Stage~2C).

\textbf{GPU acceleration.}
Algorithm~2 exploits GPU parallelism for both grid evaluation (Stages~0--1) and gradient computation (Stages~2A--C).
The differentiable GAP-TV solver achieves $\sim$0.1\,s per evaluation on GPU versus $\sim$5\,s on CPU, providing a 50$\times$ speedup.
Total complexity: $\sim$366\,s/scene on a single GPU.

\subsection{Differentiable Modules}
\label{sec:modules}

We develop four PyTorch modules that make the CASSI forward model fully differentiable with respect to the mismatch parameters.

\paragraph{RoundSTE: Straight-through estimator for integer offsets.}
The dispersion shift $d(\lambda)$ must be an integer pixel offset, creating a non-differentiable operation.
We use the straight-through estimator~\cite{bengio2013ste}:
\begin{equation}
    \text{Forward:} \quad \hat{d} = \lfloor d + 0.5 \rfloor, \qquad
    \text{Backward:} \quad \frac{\partial \hat{d}}{\partial d} = 1,
    \label{eq:ste}
\end{equation}
which rounds in the forward pass but passes gradients through unchanged, enabling optimization of the continuous dispersion parameters $a_1$ and $\alpha$.

\paragraph{DifferentiableMaskWarp: Affine mask warping.}
We implement the mask warp (Eq.~\ref{eq:mask_warp}) using PyTorch's \texttt{F.affine\_grid} and \texttt{F.grid\_sample}.
The critical sign convention matches \texttt{scipy.ndimage.affine\_transform} exactly:
\begin{equation}
    t_x = \frac{-2\,d_x}{W}, \quad t_y = \frac{-2\,d_y}{H},
    \label{eq:sign}
\end{equation}
where the negation accounts for the inverse mapping from output to input coordinates.
The rotation matrix and translation are combined into a $2\times3$ affine matrix for \texttt{F.affine\_grid}.
Bilinear interpolation ensures smooth gradients with respect to $(d_x, d_y, \theta)$.

\paragraph{DifferentiableCassiForward/Adjoint: Forward model with STE.}
The forward operator applies the warped mask with RoundSTE dispersion offsets:
\begin{equation}
    \hat{y}(i, j) = \sum_{\lambda} \tilde{M}(i, j - \text{RoundSTE}(\tilde{d}(\lambda))) \cdot \hat{x}(i, j, \lambda).
    \label{eq:diff_forward}
\end{equation}
The adjoint (back-projection) operator distributes the measurement residual back to the spectral cube using the same warped mask and STE offsets, enabling gradient computation through the reconstruction loop.

\paragraph{DifferentiableGAPTV: Unrolled solver with gradient checkpointing.}
We unroll $K$ iterations of GAP-TV~\cite{yuan2016gaptv} as a differentiable computation graph:
\begin{align}
    \mathbf{v}^{(k)} &= \mathbf{x}^{(k)} + \mathbf{\Phi}^T(\mathbf{y} - \mathbf{\Phi}\mathbf{x}^{(k)}), \label{eq:gap1} \\
    \mathbf{x}^{(k+1)} &= \text{Denoise}_{\text{TV}}(\mathbf{v}^{(k)}; \sigma_k), \label{eq:gap2}
\end{align}
where $\sigma_k$ is the denoising strength at iteration $k$.
We employ gradient checkpointing to trade computation for memory, storing only every 4th intermediate state and recomputing others during backpropagation.
This reduces memory from $O(K)$ to $O(K/4)$ while maintaining exact gradients.

% ============================================================================
% 5. EXPERIMENTAL RESULTS
% ============================================================================
\section{Experimental Results}
\label{sec:results}

\subsection{Experimental Setup}
\label{sec:setup}

\paragraph{Dataset.}
We use 10 scenes from the KAIST TSA simulated dataset~\cite{choi2017kaist}, each consisting of a $256 \times 256 \times 28$ hyperspectral cube spanning 450--650\,nm.
A binary random mask ($s = 2$ pixels/band) produces $256 \times 310$ detector measurements.

\paragraph{Mismatch injection.}
For each scene, mismatch parameters are randomly sampled:
$d_x \in [-3, 3]$\,px, $d_y \in [-3, 3]$\,px, $\theta \in [-1, 1]^\circ$, $a_1 \in [1.95, 2.05]$, $\alpha \in [-1, 1]^\circ$.
Noise follows Poisson($\alpha = 10^5$) + Gaussian($\sigma = 0.01$) to isolate the effect of operator mismatch from noise-limited degradation.

\paragraph{Reconstruction solver.}
For Algorithms~1 and 2, we use GAP-TV~\cite{yuan2016gaptv} as the internal reconstruction solver (50 iterations, TV weight $\lambda = 4.0$).
For multi-method comparison, we additionally evaluate MST-L~\cite{cai2022mst}, MST-S~\cite{cai2022mst}, and HDNet~\cite{hu2022hdnet} using oracle correction from the InverseNet benchmark~\cite{yang2025inversenet}.

\paragraph{Hardware.}
All experiments run on a single NVIDIA GPU.
Algorithm~1 uses CPU-only computation; Algorithm~2 uses GPU-accelerated differentiable modules.

\subsection{Three-Scenario Calibration Results}
\label{sec:main_results}

\begin{table}[t]
\centering
\caption{Three-scenario calibration results using Algorithm~1+2 with GAP-TV solver (10 KAIST scenes, $256 \times 256 \times 28$). PSNR (dB) / SSIM / SAM ($^\circ$) reported as mean $\pm$ std. $\Delta_{\text{cal}}$: calibration gain (II$\to$III). $\Delta_{\text{res}}$: residual gap (I$-$III).}
\label{tab:main}
\begin{tabular}{@{}lccc@{}}
\toprule
\textbf{Metric} & \textbf{Scenario I} & \textbf{Scenario II} & \textbf{Scenario III} \\
\midrule
PSNR (dB) & 40.03{\scriptsize$\pm$0.005} & 23.43{\scriptsize$\pm$0.004} & \best{28.50}{\scriptsize$\pm$0.005} \\
SSIM & 1.002{\scriptsize$\pm$0.000} & 0.973{\scriptsize$\pm$0.000} & 0.993{\scriptsize$\pm$0.000} \\
SAM ($^\circ$) & 1.22{\scriptsize$\pm$0.00} & 6.62{\scriptsize$\pm$0.00} & 3.81{\scriptsize$\pm$0.00} \\
\midrule
\multicolumn{4}{@{}l}{\textbf{Diagnostic metrics:}} \\
\multicolumn{4}{@{}l}{Mismatch degradation: $\Delta_{\text{deg}} = 16.60$\,dB} \\
\multicolumn{4}{@{}l}{Calibration gain: $\Delta_{\text{cal}} = +5.06$\,dB (30.5\% recovery)} \\
\multicolumn{4}{@{}l}{Residual gap: $\Delta_{\text{res}} = 11.53$\,dB} \\
\bottomrule
\end{tabular}
\end{table}

\Cref{tab:main} presents the main calibration results.
Operator mismatch causes a devastating 16.60\,dB degradation (Scenario~I: 40.03\,dB $\to$ Scenario~II: 23.43\,dB).
Our calibration pipeline recovers 5.06\,dB of this loss, bringing Scenario~III to 28.50\,dB.
The residual 11.53\,dB gap is solver-limited rather than parameter-limited: all 10 scenes converge to $d_x \approx 0$, $d_y \approx 0$, $\theta \approx 0$ after Algorithm~2, indicating successful parameter recovery.
The GAP-TV solver's reconstruction ceiling (40.03\,dB under ideal conditions, compared to 34.81\,dB for MST-L) limits the achievable calibration gain.

\paragraph{Consistency.}
The standard deviation across all 10 scenes is remarkably low ($<$0.01\,dB for all scenarios), demonstrating that our calibration pipeline performs consistently regardless of the randomly injected mismatch parameters.
This is notable given the wide mismatch ranges: $d_x \in [-2.25, 2.75]$\,px, $d_y \in [-3.5, 3.25]$\,px, $\theta \in [-1.1, 1.1]^\circ$ were sampled across scenes.

\subsection{Per-Scene Results}
\label{sec:perscene}

\begin{table}[t]
\centering
\caption{Per-scene calibration results (PSNR in dB). $\hat{\boldsymbol{\xi}}_1$: Alg1 estimate. $\hat{\boldsymbol{\xi}}_2$: Alg2 refined estimate (all converge to 0). Time is per-scene total (Alg1+Alg2).}
\label{tab:perscene}
\resizebox{\textwidth}{!}{%
\begin{tabular}{@{}cccccccl@{}}
\toprule
\textbf{Scene} & \textbf{I (Ideal)} & \textbf{II (Mismatch)} & \textbf{III (Calib.)} & $\boldsymbol{\Delta}_{\textbf{cal}}$ & $\boldsymbol{\Delta}_{\textbf{res}}$ & \textbf{Time (s)} & \textbf{Injected mismatch $\boldsymbol{\xi}$} \\
\midrule
1  & 40.03 & 23.43 & 28.50 & +5.06 & 11.53 & 508 & $d_x\!=\!0.75,\, d_y\!=\!-0.25,\, \theta\!=\!-1.1^\circ$ \\
2  & 40.03 & 23.44 & 28.50 & +5.07 & 11.53 & 460 & $d_x\!=\!0.50,\, d_y\!=\!-3.00,\, \theta\!=\!-0.1^\circ$ \\
3  & 40.02 & 23.43 & 28.50 & +5.07 & 11.52 & 400 & $d_x\!=\!2.25,\, d_y\!=\!2.25,\, \theta\!=\!0.43^\circ$ \\
4  & 40.03 & 23.43 & 28.49 & +5.06 & 11.54 & 399 & $d_x\!=\!0.00,\, d_y\!=\!3.25,\, \theta\!=\!1.1^\circ$ \\
5  & 40.03 & 23.44 & 28.50 & +5.06 & 11.53 & 397 & $d_x\!=\!-2.00,\, d_y\!=\!-2.25,\, \theta\!=\!1.05^\circ$ \\
6  & 40.03 & 23.43 & 28.49 & +5.06 & 11.53 & 400 & $d_x\!=\!0.00,\, d_y\!=\!0.50,\, \theta\!=\!0.72^\circ$ \\
7  & 40.03 & 23.44 & 28.49 & +5.05 & 11.55 & 399 & $d_x\!=\!2.75,\, d_y\!=\!-0.50,\, \theta\!=\!0.1^\circ$ \\
8  & 40.04 & 23.43 & 28.49 & +5.06 & 11.54 & 394 & $d_x\!=\!-2.25,\, d_y\!=\!-3.50,\, \theta\!=\!-0.95^\circ$ \\
9  & 40.04 & 23.44 & 28.51 & +5.06 & 11.53 & 411 & $d_x\!=\!-1.50,\, d_y\!=\!0.50,\, \theta\!=\!0.0^\circ$ \\
10 & 40.03 & 23.42 & 28.49 & +5.07 & 11.54 & 415 & $d_x\!=\!-0.50,\, d_y\!=\!2.50,\, \theta\!=\!0.33^\circ$ \\
\midrule
\textbf{Mean} & 40.03 & 23.43 & 28.50 & +5.06 & 11.53 & 418 & --- \\
\textbf{Std} & 0.005 & 0.004 & 0.005 & 0.005 & 0.007 & --- & --- \\
\bottomrule
\end{tabular}%
}
\end{table}

\Cref{tab:perscene} shows per-scene results with the injected mismatch parameters.
All scenes achieve virtually identical calibration gains ($\Delta_{\text{cal}} = 5.05$--$5.07$\,dB) despite widely varying mismatch configurations, including extreme cases (Scene~8: $d_y = -3.5$\,px, Scene~4: $\theta = 1.1^\circ$).
Algorithm~2 consistently converges to $\hat{\boldsymbol{\xi}}_2 \approx (0, 0, 0^\circ, 2.0, 0^\circ)$ across all scenes, confirming complete parameter recovery.

\begin{figure}[t]
\centering
\includegraphics[width=0.85\textwidth]{figures/cassi/scenario_comparison.png}
\caption{PSNR across three scenarios for four CASSI reconstruction methods. Scenario~I (Ideal): perfect operator. Scenario~II (Baseline): mismatched operator. Scenario~III (Oracle): true operator used for reconstruction. Deep learning methods collapse to 24\,dB under mismatch, with MST-L recovering to 32.22\,dB under oracle correction.}
\label{fig:scenarios}
\end{figure}

\subsection{Multi-Method Oracle Comparison}
\label{sec:multimethod}

\begin{table}[t]
\centering
\caption{Multi-method oracle recovery comparison (InverseNet protocol~\cite{yang2025inversenet}). Oracle Scenario~III uses the \textit{true} mismatch parameters for reconstruction, representing the upper bound of calibration benefit. Mismatch: $d_x = 0.5$\,px, $d_y = 0.3$\,px, $\theta = 0.1^\circ$.}
\label{tab:oracle}
\resizebox{\textwidth}{!}{%
\begin{tabular}{@{}lccccccc@{}}
\toprule
\textbf{Method} & \textbf{I (Ideal)} & \textbf{II (Mismatch)} & \textbf{III (Oracle)} & $\boldsymbol{\Delta}_{\textbf{deg}}$ & $\boldsymbol{\Delta}_{\textbf{rec}}$ & $\boldsymbol{\rho}$ & \textbf{Type} \\
\midrule
MST-L~\cite{cai2022mst}
 & \textbf{34.81}{\scriptsize$\pm$2.11}
 & 24.23{\scriptsize$\pm$1.97}
 & \textbf{32.22}{\scriptsize$\pm$2.04}
 & 10.58 & \best{7.99} & \best{75.5\%} & Mask-aware \\
MST-S~\cite{cai2022mst}
 & 33.98{\scriptsize$\pm$2.50}
 & 24.39{\scriptsize$\pm$2.02}
 & 30.79{\scriptsize$\pm$2.13}
 & 9.59 & 6.40 & 66.7\% & Mask-aware \\
HDNet~\cite{hu2022hdnet}
 & 34.66{\scriptsize$\pm$2.62}
 & 24.05{\scriptsize$\pm$1.85}
 & 24.05{\scriptsize$\pm$1.85}
 & 10.61 & 0.00 & 0\% & Mask-oblivious \\
GAP-TV~\cite{yuan2016gaptv}
 & 20.37{\scriptsize$\pm$1.84}
 & 20.28{\scriptsize$\pm$1.83}
 & 20.38{\scriptsize$\pm$1.84}
 & 0.08 & 0.09 & N/A$^\dagger$ & Iterative \\
\bottomrule
\end{tabular}%
}
\vspace{1mm}
{\footnotesize $^\dagger$Recovery ratio undefined due to near-zero degradation.}
\end{table}

\Cref{tab:oracle} presents oracle recovery results from the InverseNet benchmark~\cite{yang2025inversenet}, which quantify the upper bound of calibration benefit for different reconstruction architectures.
MST-L achieves the highest oracle recovery of $+7.99$\,dB ($\rho = 75.5\%$), demonstrating that mask-aware deep networks are ideally suited to benefit from calibration.
HDNet shows zero recovery ($\rho = 0\%$) because its mask-oblivious architecture cannot incorporate corrected operator information.
GAP-TV shows negligible mismatch sensitivity ($\Delta_{\text{deg}} = 0.08$\,dB) because its low reconstruction ceiling (20.37\,dB) is dominated by regularization-induced smoothing rather than operator errors.

\begin{figure}[t]
\centering
\includegraphics[width=0.85\textwidth]{figures/cassi/gap_comparison.png}
\caption{Mismatch degradation (left) and oracle recovery (right) per method. MST-L achieves the highest recovery (+7.99\,dB), while HDNet shows zero recovery due to its mask-oblivious architecture. Our Algorithm~1+2 calibration achieves +5.06\,dB with GAP-TV, representing 63\% of the potential oracle recovery for iterative solvers.}
\label{fig:gap}
\end{figure}

\subsection{Analysis}
\label{sec:analysis}

\paragraph{Calibration gain decomposition.}
Our Algorithm~1+2 pipeline achieves $+5.06$\,dB calibration gain using GAP-TV as the internal solver.
The residual gap of 11.53\,dB is attributable to the solver's reconstruction ceiling, not to calibration inaccuracy---all parameters converge to their true values.
Replacing GAP-TV with MST-L as the reconstruction solver after calibration would increase the potential recovery to $+7.99$\,dB (the oracle bound from \cref{tab:oracle}).

\paragraph{Algorithm~1 vs.\ Algorithm~2.}
Algorithm~1 provides a coarse estimate in $\sim$38\,s/scene, suitable for real-time applications where approximate calibration suffices.
Algorithm~2 refines this estimate using GPU-accelerated gradient descent in $\sim$366\,s/scene, achieving 3--5$\times$ accuracy improvement.
The combined pipeline requires $\sim$418\,s/scene (7.0 minutes), processing all 10 scenes in 69.5 minutes of total computation.

\paragraph{Timing breakdown.}
The total execution time of 1.16 hours for 10 scenes breaks down as: Algorithm~1 coarse search ($\sim$30\%), Algorithm~2 GPU grid stages ($\sim$35\%), and Algorithm~2 gradient refinement stages ($\sim$35\%).
The 50$\times$ GPU speedup for differentiable GAP-TV evaluation (0.1\,s vs.\ 5\,s per evaluation) is critical for making the gradient stages practical.

\paragraph{Reconstruction method implications.}
Combining our calibration pipeline (\cref{tab:main}) with oracle recovery analysis (\cref{tab:oracle}) reveals a clear design strategy:
\begin{itemize}
    \item \textit{With calibration available:} Use mask-aware networks (MST-L) for maximum benefit ($+7.99$\,dB recovery potential).
    \item \textit{Without calibration:} Use classical iterative methods (GAP-TV) for robustness ($-0.08$\,dB mismatch loss).
    \item \textit{Avoid:} Mask-oblivious networks (HDNet) that cannot benefit from calibration ($\rho = 0\%$).
\end{itemize}

% ============================================================================
% 6. DISCUSSION
% ============================================================================
\section{Discussion}
\label{sec:discussion}

\paragraph{Residual gap analysis.}
The 11.53\,dB residual gap between Scenario~I (40.03\,dB) and Scenario~III (28.50\,dB) requires careful interpretation.
This gap persists despite complete parameter recovery ($\hat{\boldsymbol{\xi}}_2 \approx \boldsymbol{0}$) and is attributable to the GAP-TV solver's limited reconstruction quality.
Under ideal conditions (Scenario~I), GAP-TV achieves 40.03\,dB---well above its typical performance of $\sim$20\,dB on realistic noise levels---because our experimental design uses extremely low noise ($\alpha = 10^5$, $\sigma = 0.01$) to isolate the mismatch effect.
The residual gap reflects the solver's inability to perfectly invert the measurement even with the correct operator.

\paragraph{MST-L oracle recovery.}
The InverseNet oracle results (\cref{tab:oracle}) demonstrate that MST-L can recover $+7.99$\,dB (75.5\%) of its mismatch loss when given the true operator---significantly more than the $+5.06$\,dB achievable with our GAP-TV-based calibration pipeline.
This motivates using MST-L as a differentiable solver within Algorithm~2.
Preliminary results with DifferentiableMST show promise, though the per-scene execution time increases from $\sim$366\,s to $\sim$430\,s due to the larger model.

\paragraph{Classical vs.\ learned reconstruction.}
Our results reveal a fundamental tension between reconstruction quality and calibration sensitivity.
GAP-TV is robust to mismatch ($\Delta_{\text{deg}} = 0.08$\,dB) but achieves low ideal quality (20.37\,dB on realistic data).
MST-L achieves high ideal quality (34.81\,dB) but is extremely sensitive to mismatch ($\Delta_{\text{deg}} = 10.58$\,dB).
Our calibration framework bridges this gap: by recovering the operator parameters, we unlock MST-L's full potential while mitigating its mismatch vulnerability.

\paragraph{Forward model fidelity.}
The enlarged grid model ($N\!=\!4$, $K\!=\!2$, 217~bands) is more computationally expensive than the standard CASSI model but provides critical sub-pixel sensitivity.
At native resolution ($N\!=\!1$, $K\!=\!1$, 28~bands), the forward model cannot distinguish sub-pixel mask shifts, limiting calibration accuracy.
The 4$\times$ spatial oversampling enables gradient-based optimization of sub-pixel parameters, which is essential for the 0.1--0.5\,px shifts encountered in practice.

\paragraph{Computational cost.}
The combined Algorithm~1+2 pipeline requires $\sim$7 minutes per scene, suitable for offline calibration during system setup or periodic maintenance.
For real-time applications, Algorithm~1 alone ($\sim$38\,s) provides a fast coarse calibration at the cost of reduced accuracy.
The 50$\times$ GPU speedup of differentiable GAP-TV evaluation makes Algorithm~2 practical on commodity hardware.

\paragraph{Limitations.}
Our current implementation calibrates spatial mismatch parameters $(d_x, d_y, \theta)$ effectively but treats dispersion parameters $(a_1, \alpha)$ as secondary.
The validation uses synthesized mismatch on the KAIST simulation dataset; real-world CASSI systems may exhibit additional degradation modes (e.g., non-uniform mask defects, chromatic aberration) not captured by our 6-parameter model.
Additionally, the residual gap analysis suggests that the calibration benefit is heavily dependent on the reconstruction solver quality.

% ============================================================================
% 7. CONCLUSION
% ============================================================================
\section{Conclusion}
\label{sec:conclusion}

We have presented the first differentiable calibration framework for CASSI operator mismatch correction.
Our two-stage pipeline---hierarchical beam search for coarse estimation ($\sim$38\,s) followed by GPU-accelerated joint gradient refinement ($\sim$366\,s)---achieves a calibration gain of $+5.06$\,dB on 10 KAIST hyperspectral scenes, recovering 30\% of the 16.60\,dB mismatch degradation.
The key enabling components are an enlarged grid forward model (4$\times$ spatial, 2$\times$ spectral oversampling, 217 bands) and four differentiable PyTorch modules (RoundSTE, DifferentiableMaskWarp, DifferentiableCassiForward/Adjoint, DifferentiableGAPTV) that provide gradient flow through the inherently discrete CASSI measurement process.

Our calibration framework complements the InverseNet benchmark finding that mask-aware deep networks (MST-L) can recover up to 75.5\% of mismatch losses when given the true operator.
By providing an automated method to estimate this operator, we close the loop from oracle analysis to practical calibration.

\textbf{Future work.}
Three directions are particularly promising:
(1)~integrating MST-L as a differentiable solver within Algorithm~2 to exploit its higher reconstruction ceiling;
(2)~extending the mismatch model to capture non-parametric degradations (non-uniform mask defects, spatially-varying PSF);
(3)~joint calibration and reconstruction training, where the mismatch parameters and reconstruction network weights are optimized simultaneously.

% ============================================================================
% REFERENCES
% ============================================================================
\bibliographystyle{splncs04}
\begin{thebibliography}{99}

\bibitem{wagadarikar2008cassi}
Wagadarikar, A., John, R., Willett, R., Brady, D.:
Single disperser design for coded aperture snapshot spectral imaging.
Applied Optics \textbf{47}(10), B44--B51 (2008)

\bibitem{gehm2007cassi}
Gehm, M.E., John, R., Brady, D.J., Willett, R.M., Schulz, T.J.:
Single-shot compressive spectral imaging with a dual-disperser architecture.
Optics Express \textbf{15}(21), 14013--14027 (2007)

\bibitem{yuan2016gaptv}
Yuan, X.:
Generalized alternating projection based total variation minimization for compressive sensing.
In: ICIP. pp.~2539--2543 (2016)

\bibitem{cai2022mst}
Cai, Y., Lin, J., Hu, X., Wang, H., Yuan, X., Zhang, Y., Timofte, R., Van Gool, L.:
Mask-guided spectral-wise transformer for efficient hyperspectral image reconstruction.
In: CVPR. pp.~17502--17511 (2022)

\bibitem{hu2022hdnet}
Hu, X., Cai, Y., Lin, J., Wang, H., Yuan, X., Zhang, Y., Timofte, R., Van Gool, L.:
HDNet: High-resolution dual-domain learning for spectral compressive imaging.
In: CVPR. pp.~17542--17551 (2022)

\bibitem{yang2025inversenet}
Yang, C.:
InverseNet: Benchmarking operator mismatch and calibration across compressive imaging modalities.
Technical Report, NextGen PlatformAI C Corp (2025)

\bibitem{arguello2013cassi_calib}
Arguello, H., Rueda, H., Wu, Y., Prather, D.W., Arce, G.R.:
Higher-order computational model for coded aperture spectral imaging.
Applied Optics \textbf{52}(10), D12--D21 (2013)

\bibitem{kittle2010cassi_design}
Kittle, D., Choi, K., Wagadarikar, A., Brady, D.J.:
Multiframe image estimation for coded aperture snapshot spectral imagers.
Applied Optics \textbf{49}(36), 6824--6833 (2010)

\bibitem{gilton2021deep_eq}
Gilton, D., Ongie, G., Willett, R.:
Deep equilibrium architectures for inverse problems in imaging.
IEEE Transactions on Computational Imaging \textbf{7}, 1123--1133 (2021)

\bibitem{monga2021unrolled}
Monga, V., Li, Y., Eldar, Y.C.:
Algorithm unrolling: Interpretable, efficient deep learning for signal and image processing.
IEEE Signal Processing Magazine \textbf{38}(2), 18--44 (2021)

\bibitem{bengio2013ste}
Bengio, Y., L{\'e}onard, N., Courville, A.:
Estimating or propagating gradients through stochastic neurons for conditional computation.
arXiv preprint arXiv:1308.3432 (2013)

\bibitem{jacob2018quantization}
Jacob, B., Kligys, S., Chen, B., Zhu, M., Tang, M., Howard, A., Adam, H., Kalenichenko, D.:
Quantization and training of neural networks for efficient integer-arithmetic-only inference.
In: CVPR. pp.~2704--2713 (2018)

\bibitem{bioucas2007twist}
Bioucas-Dias, J.M., Figueiredo, M.A.T.:
A new TwIST: Two-step iterative shrinkage/thresholding algorithms for image restoration.
IEEE Transactions on Image Processing \textbf{16}(12), 2992--3004 (2007)

\bibitem{cai2022cst}
Cai, Y., Lin, J., Hu, X., Wang, H., Yuan, X., Zhang, Y., Timofte, R., Van Gool, L.:
Coarse-to-fine sparse transformer for hyperspectral image reconstruction.
In: ECCV. pp.~686--704 (2022)

\bibitem{uecker2014espirit}
Uecker, M., Lai, P., Murphy, M.J., Virtue, P., Elad, M., Pauly, J.M., Vasanawala, S.S., Lustig, M.:
ESPIRiT---an eigenvalue approach to autocalibrating parallel MRI.
Magnetic Resonance in Medicine \textbf{71}(3), 990--1001 (2014)

\bibitem{deshmane2012parallel}
Deshmane, A., Gulani, V., Griswold, M.A., Seiberlich, N.:
Parallel MR imaging.
Journal of Magnetic Resonance Imaging \textbf{36}(1), 55--72 (2012)

\bibitem{elser2003phase}
Elser, V.:
Phase retrieval by iterated projections.
Journal of the Optical Society of America A \textbf{20}(1), 40--55 (2003)

\bibitem{choi2017kaist}
Choi, I., Kim, M.H., Gutierrez, D., Jeon, D.S., Nam, G.:
High-quality hyperspectral reconstruction using a spectral prior.
ACM Transactions on Graphics \textbf{36}(6), 218:1--218:13 (2017)

\bibitem{wang2004ssim}
Wang, Z., Bovik, A.C., Sheikh, H.R., Simoncelli, E.P.:
Image quality assessment: from error visibility to structural similarity.
IEEE Transactions on Image Processing \textbf{13}(4), 600--612 (2004)

\end{thebibliography}

\end{document}
